\documentclass[11pt,a4paper]{article}

% ─── Packages ──────────────────────────────────────────────────────────────────
\usepackage[utf8]{inputenc}
\usepackage[T1]{fontenc}
\usepackage{amsmath,amssymb}
\usepackage{graphicx}
\usepackage{booktabs}
\usepackage{hyperref}
\usepackage{natbib}
\usepackage[margin=1in]{geometry}
\usepackage{algorithm}
\usepackage{algpseudocode}
\usepackage{xcolor}
\usepackage{enumitem}
\usepackage{microtype}
\usepackage{pifont}
\usepackage{multirow}

\hypersetup{
    colorlinks=true,
    linkcolor=blue,
    citecolor=blue,
    urlcolor=blue,
}

% ─── Title ─────────────────────────────────────────────────────────────────────

\title{The Retrieval--Reasoning Gap: Why Better Retrieval Does Not Guarantee\\Better Reasoning in Iterative RAG Systems\thanks{Code and data: \url{https://github.com/yakkalaanugna/Automated-RAG-Agent}}}

\author{Anugna Yakkala \\[4pt]
{\normalsize School of Computer Science and Engineering,}\\
{\normalsize Vellore Institute of Technology, Vellore, TN, India}}

\date{April 2026}

\begin{document}
\maketitle

% ═══════════════════════════════════════════════════════════════════════════════
% ABSTRACT
% ═══════════════════════════════════════════════════════════════════════════════

\begin{abstract}
Retrieval-augmented generation (RAG) evaluations implicitly assume that
improving retrieval improves end-task accuracy.  We show this assumption
breaks down for multi-hop RAG.
Using an iterative pipeline (hybrid BM25\,+\,bi-encoder retrieval,
cross-encoder reranking, confidence-gated iteration), we evaluate
65 multi-hop root-cause queries on a synthetic telecom log corpus
(1{,}115 entries, 10 failure scenarios).
A six-configuration ablation with a fixed LLM reveals that
Precision@K spans 0.553--0.794 ($\Delta{=}0.24$,
$p{<}0.001$) while root-cause accuracy remains flat at
0.311--0.359 ($\Delta{=}0.05$, $p{=}0.08$, n.s.).
A multi-LLM comparison (8B--70B) with identical retrieval confirms
that reasoning accuracy depends on model capability independently
of evidence quality.
We propose a composite evaluation metric (semantic similarity +
ROUGE-L + structured scoring) that better captures diagnostic
quality than keyword overlap alone ($r{=}0.67$).
Our results establish the \emph{retrieval--reasoning gap} as a
structural property of multi-hop RAG.
\end{abstract}

% ═══════════════════════════════════════════════════════════════════════════════
% 1. INTRODUCTION
% ═══════════════════════════════════════════════════════════════════════════════

\section{Introduction}

Evaluation of retrieval-augmented generation (RAG) systems
\citep{lewis2020rag,gao2023ragsurvey} typically reports two metric
families side by side: retrieval quality (Precision@K, Recall@K, nDCG)
and end-task quality (answer accuracy, F1).
The implicit assumption is that improving retrieval improves the
end-task.
For multi-hop RAG---where evidence must be assembled from multiple
sources before a language model reasons over it---this assumption is
load-bearing yet under-examined.

This paper provides controlled empirical evidence that the assumption
is wrong.
We present the \textbf{retrieval--reasoning gap}: the observation that
improvements in retrieval-side metrics do not translate to improvements
in reasoning accuracy when the LLM reader is fixed, and that this
gap is a structural property of multi-hop RAG rather than an artifact
of dataset size.

Our experimental design isolates retrieval from reasoning through
three mechanisms:
\begin{enumerate}[nosep]
    \item A \textbf{six-configuration ablation study} that varies the
    retrieval pipeline ($\pm$\,BM25, $\pm$\,cross-encoder reranking,
    $\pm$\,iterative refinement) while holding the LLM constant
    (Section~\ref{sec:ablation}).
    \item A \textbf{multi-LLM comparison} that fixes retrieval and
    varies only the reasoning model across four architectures spanning
    8B--70B parameters (Section~\ref{sec:multi-llm}).
    \item \textbf{Improved evaluation metrics}---semantic similarity,
    structured scoring, composite metric---that capture reasoning
    quality beyond keyword overlap (Section~\ref{sec:metrics-result}).
\end{enumerate}

We evaluate on a synthetic telecom log corpus (1{,}115 records, 10
log files, 10 failure scenarios) with 65 evaluation queries requiring
multi-hop cross-file reasoning.
The telecom setting is a controlled testbed; the evaluation phenomena
we surface---metric disagreement, failure-mode independence, and
inadequacy of aggregate scores---generalize to any RAG system
performing multi-hop evidence assembly followed by LLM reasoning.

\paragraph{Contributions.}
\begin{enumerate}[nosep]
    \item \textbf{Empirical evidence of the retrieval--reasoning gap.}
    Across six ablation configurations, Precision@K varies by 0.24
    ($p{<}0.001$) while root-cause accuracy varies by only 0.05
    ($p{=}0.08$, n.s.; Pearson $r{=}0.244$).
    Multi-LLM comparison with fixed retrieval shows reasoning
    accuracy gaps of up to 0.103 attributable to model capability
    alone.
    \item \textbf{Evaluation framework.}
    We propose a decomposition protocol for multi-hop RAG evaluation:
    (a)~ablation with fixed LLM, (b)~model comparison with fixed
    retrieval, (c)~per-query regime classification (high/low
    retrieval $\times$ correct/incorrect reasoning).
    \item \textbf{Composite evaluation metric.}
    We show keyword overlap is a noisy proxy and propose a composite
    score (semantic similarity + ROUGE-L + structured scoring) that
    better captures diagnostic quality.
    \item \textbf{Dataset release.}
    Synthetic telecom log dataset: 1{,}115 entries, 10 failure
    scenarios, 65 annotated queries with ground-truth root causes.
\end{enumerate}

% ═══════════════════════════════════════════════════════════════════════════════
% 2. RELATED WORK
% ═══════════════════════════════════════════════════════════════════════════════

\section{Related Work}
\label{sec:related}

\paragraph{Retrieval-Augmented Generation.}
\citet{lewis2020rag} showed that augmenting LMs with retrieved
passages yields large gains on knowledge-intensive tasks;
\citet{gao2023ragsurvey} survey extensions to summarization, code
generation, and domain-specific applications.
Most RAG systems follow a single-pass architecture insufficient for
multi-file log analysis.

\paragraph{Iterative and Multi-Step Retrieval.}
Self-RAG \citep{asai2024selfrag} enables self-reflective retrieval;
Active RAG \citep{jiang2023activerar} uses active learning signals;
query rewriting \citep{wang2023queryrewrite} reformulates queries
after initial passes.
These approaches lack principled stopping criteria.
Our adaptive confidence-gated iteration addresses this, and our
ablation quantifies the contribution of iteration to both retrieval
and reasoning.

\paragraph{Hybrid Retrieval and Cross-Encoder Reranking.}
Fusing BM25 \citep{robertson2009bm25} and bi-encoder retrieval
\citep{reimers2019sentencebert} is standard practice.
Cross-encoder rerankers \citep{nogueira2019passage} substantially
improve precision; BEIR \citep{thakur2021beir} confirms zero-shot
generalization.
We show that reranking improves \emph{retrieval} precision but does
not proportionally improve \emph{reasoning} accuracy.

\paragraph{RAG Evaluation.}
RAGAS \citep{es2024ragas} and ARES \citep{saad2024ares} decompose
RAG quality into context-relevance, faithfulness, and answer-relevance.
Lost-in-the-middle \citep{liu2024lostmiddle} shows LLMs underuse
mid-context evidence---invisible to Precision@K.
We complement this literature by demonstrating that retrieval and
end-task metrics move independently, and that neither alone nor
together reveals which failure mode is responsible.

% ═══════════════════════════════════════════════════════════════════════════════
% 3. METHOD
% ═══════════════════════════════════════════════════════════════════════════════

\section{Method}
\label{sec:method}

\subsection{Pipeline Architecture}

The system implements a four-stage iterative loop:
\emph{Hybrid Retrieval $\to$ Cross-Encoder Reranking $\to$ LLM
Reasoning $\to$ Confidence-Gated Iteration}.

\paragraph{Stage 1: Hybrid Candidate Generation.}
Dense retrieval (\texttt{all-MiniLM-L6-v2}, 384-dim) indexed in
ChromaDB fuses with BM25 ($k_1{=}1.5$, $b{=}0.75$):
$s_{\text{hybrid}}(d) = 0.7 \cdot s_{\text{dense}}(d) + 0.3 \cdot
s_{\text{BM25}}(d)$.
Over-fetches $N{=}20$--$30$ candidates for reranking.

\paragraph{Stage 2: Cross-Encoder Reranking.}
\texttt{ms-marco-MiniLM-L-6-v2} (22M params) rescores each
(query, document) pair via joint cross-attention.
The cross-encoder score is the \emph{sole} ranking signal,
discarding hybrid scores to avoid reintroducing first-stage noise.
Top-$k{=}6$ documents proceed to the LLM.

\paragraph{Stage 3: LLM Analysis.}
The LLM (default: Llama~3.3~70B, temperature~0) produces structured
output: root cause, severity, timeline, reasoning steps,
recommendation.
Composite confidence combines: mean retrieval score (30\%), output
completeness (30\%), evidence density (20\%), inter-iteration
consistency (20\%).

\paragraph{Stage 4: Confidence-Gated Iteration.}
Three stopping rules: high-confidence exit (${\geq}0.85$), plateau
detection (2 consecutive non-improvements), hard cap ($M{=}3$).
Returns the \emph{best-scoring} iteration.

\begin{algorithm}[t]
\caption{Adaptive Iterative RAG}
\label{alg:adaptive}
\begin{algorithmic}[1]
\Require Query $q$, documents $\mathcal{D}$, max iters $M$, threshold $\epsilon$
\State $q_1 \gets q$; $c_0 \gets 0$; $\text{best} \gets \textsc{null}$; $\text{seen} \gets \emptyset$
\For{$i = 1$ \textbf{to} $M$}
    \State $\mathcal{C}_i \gets \textsc{HybridRetrieve}(q_i, \mathcal{D}, N_i)$
    \State $\mathcal{R}_i \gets \textsc{Rerank}(q_i, \mathcal{C}_i, k)$
    \State $\mathcal{R}_i \gets \textsc{Dedup}(\mathcal{R}_i, \text{seen})$; update $\text{seen}$
    \State $a_i \gets \textsc{LLM}(q_i, \mathcal{R}_i)$
    \State $c_i \gets \textsc{Confidence}(a_i, \mathcal{R}_i, a_{i-1})$
    \If{$c_i > c_{\text{best}}$} $\text{best} \gets (a_i, \mathcal{R}_i, c_i)$ \EndIf
    \If{$c_i \geq 0.85$ \textbf{or} drops $\geq 2$} \textbf{break} \EndIf
    \State $q_{i+1} \gets \textsc{RefineQuery}(q, q_i, a_i)$
\EndFor
\State \Return $\text{best}$
\end{algorithmic}
\end{algorithm}

% ═══════════════════════════════════════════════════════════════════════════════
% 4. EXPERIMENTAL SETUP
% ═══════════════════════════════════════════════════════════════════════════════

\section{Experimental Setup}
\label{sec:setup}

\subsection{Dataset}

We construct a synthetic telecom log corpus for multi-hop retrieval
evaluation: \textbf{1{,}115 log entries} across \textbf{10 log files}
from six network domains (eNodeB, CU-CP, CU-UP, RAN, Core Network,
Transport).
The corpus encodes \textbf{10 failure scenarios}, each comprising a
causal chain spanning 2--5 files:

\begin{enumerate}[nosep]
    \item RRC Reconfiguration failure cascade (4 files, 13 steps)
    \item SCTP transport link failure, 45 UEs affected (4 files)
    \item Inter-gNB handover failure with PDCP data loss (4 files)
    \item PDU session rejection by PCF policy (5 files)
    \item RLC max retransmission causing DRB failure (4 files)
    \item GTP-U path failure on backhaul (4 files)
    \item 5G-AKA authentication failure, SQN desync (4 files)
    \item DL scheduler starvation under load (3 files)
    \item IPSec tunnel failure on midhaul (3 files)
    \item AMF overload from registration storm (3 files)
\end{enumerate}

Each file contains scenario entries interleaved with 80--120
normal-operation background entries, requiring retrieval to
distinguish relevant evidence from noise.

\subsection{Queries}

\textbf{65 evaluation queries} across 8 types:
root-cause analysis (10), failure tracing (10), multi-hop (10),
temporal (10), impact analysis (10), error-code (10),
cross-scenario (4), contrastive (1).
55 of 65 require cross-file reasoning.
Each query is annotated with relevant files, root-cause keywords,
and free-text ground truth.

\subsection{Ablation Configurations}

Six configurations share the same LLM (Llama~3.3~70B):

\begin{table}[h]
\centering
\small
\caption{Ablation configurations.}
\label{tab:configs}
\begin{tabular}{lcccc}
\toprule
\textbf{Config} & \textbf{Dense} & \textbf{BM25} & \textbf{Rerank} & \textbf{Iter.} \\
\midrule
Dense-Only       & \checkmark &            &             &             \\
BM25-Only        &            & \checkmark &             &             \\
Hybrid           & \checkmark & \checkmark &             &             \\
Hybrid+Rerank    & \checkmark & \checkmark & \checkmark  &             \\
Hybrid+Iter.     & \checkmark & \checkmark &             & \checkmark  \\
Full System      & \checkmark & \checkmark & \checkmark  & \checkmark  \\
\bottomrule
\end{tabular}
\end{table}

\subsection{Multi-LLM Comparison}

Retrieval is fixed (Hybrid+Rerank, top-6) across four models:
Llama~3.3~70B, Llama~3.1~8B, Mixtral~8x7B, Gemma2~9B (all via Groq,
temperature~0).

\subsection{Metrics}

\textbf{Retrieval}: Precision@K, Recall@K, MRR.
\textbf{Reasoning}: keyword overlap (RCA\textsubscript{kw}),
semantic similarity, structured scoring, composite score
($0.35{\cdot}\text{sem} + 0.30{\cdot}\text{struct} +
0.20{\cdot}\text{ROUGE-L} + 0.15{\cdot}\text{kw}$).
\textbf{LLM-as-judge}: binary correctness.

% ═══════════════════════════════════════════════════════════════════════════════
% 5. RESULTS
% ═══════════════════════════════════════════════════════════════════════════════

\section{Results}
\label{sec:results}

\subsection{Ablation Study}
\label{sec:ablation}

Table~\ref{tab:ablation} reports aggregate metrics.
The central finding: \textbf{retrieval metrics move; reasoning
metrics do not}.

\begin{table}[h]
\centering
\small
\caption{Ablation results (65 queries, Llama 3.3 70B).
Values produced by \texttt{run\_ablation.py} and filled at run time.}
\label{tab:ablation}
\begin{tabular}{lcccccc}
\toprule
\textbf{Config} & \textbf{P@K} & \textbf{R@K} & \textbf{MRR} & \textbf{RCA\textsubscript{kw}} & \textbf{Conf.} & \textbf{Iters} \\
\midrule
Dense-Only           & 0.553$\pm$0.18 & 0.500$\pm$0.19 & 0.613$\pm$0.18 & 0.311$\pm$0.16 & 0.553 & 1.0 \\
BM25-Only            & 0.598$\pm$0.18 & 0.509$\pm$0.21 & 0.629$\pm$0.17 & 0.314$\pm$0.16 & 0.550 & 1.0 \\
Hybrid               & 0.699$\pm$0.16 & 0.637$\pm$0.19 & 0.757$\pm$0.14 & 0.312$\pm$0.17 & 0.588 & 1.0 \\
Hybrid+Rerank        & 0.794$\pm$0.15 & 0.720$\pm$0.17 & 0.822$\pm$0.14 & 0.344$\pm$0.16 & 0.641 & 1.0 \\
Hybrid+Iter.         & 0.628$\pm$0.19 & 0.704$\pm$0.16 & 0.716$\pm$0.16 & 0.338$\pm$0.15 & 0.606 & 2.0 \\
Full-System          & 0.751$\pm$0.15 & 0.729$\pm$0.15 & 0.811$\pm$0.14 & 0.359$\pm$0.18 & 0.649 & 2.2 \\
\bottomrule
\end{tabular}
\end{table}

\paragraph{Findings.}
(i)~Cross-encoder reranking yields the largest single-component
gain: Hybrid$\to$Hybrid+Rerank improves P@K by +0.096
($p{<}0.001$), yet RCA\textsubscript{kw} increases by only
+0.032 ($p{=}0.234$, n.s.).
(ii)~BM25 fusion adds complementary signal for exact error codes,
lifting P@K from 0.553 (Dense-Only) to 0.699 (Hybrid).
(iii)~Iteration without reranking degrades precision via query drift
(Hybrid+Iter.\ P@K$=$0.628 vs.\ Hybrid 0.699), though recall
improves (0.704 vs.\ 0.637).
(iv)~Across all six configs, \textbf{P@K varies by 0.24
($p{<}0.001$) while RCA\textsubscript{kw} varies by only 0.05
($p{=}0.08$, n.s.)}---the direct empirical signature of the
retrieval--reasoning gap.

\subsection{Multi-LLM Comparison}
\label{sec:multi-llm}

\begin{table}[h]
\centering
\small
\caption{Multi-LLM results with identical retrieval.
Values from \texttt{run\_multi\_llm\_comparison.py}.}
\label{tab:multi-llm}
\begin{tabular}{lcccc}
\toprule
\textbf{Model} & \textbf{RCA\textsubscript{kw}} & \textbf{Completeness} & \textbf{Latency} & \textbf{Params} \\
\midrule
Llama-3.3-70B        & 0.363$\pm$0.14 & 0.703$\pm$0.10 & 4.4 & 70B \\
Mixtral-8x7B         & 0.300$\pm$0.12 & 0.674$\pm$0.09 & 3.8 & 46.7B \\
Gemma2-9B            & 0.302$\pm$0.13 & 0.632$\pm$0.10 & 2.1 & 9B \\
Llama-3.1-8B         & 0.260$\pm$0.12 & 0.607$\pm$0.12 & 1.9 & 8B \\
\bottomrule
\end{tabular}
\end{table}

With identical Hybrid+Rerank retrieval (P@K$\approx$0.794 for all
models), RCA\textsubscript{kw} ranges from 0.260 (Llama~3.1~8B) to
0.363 (Llama~3.3~70B)---a 0.103 absolute gap attributable entirely
to model capability.
Completeness follows the same trend (0.607--0.703), confirming that
\textbf{reasoning is independently bounded by model capability}.
Notably, the 70B$\to$8B reasoning drop (0.103) exceeds the
Dense-Only$\to$Hybrid+Rerank retrieval-driven RCA change (0.032),
indicating that model scaling dominates retrieval optimization for
reasoning quality.

\subsection{Retrieval--Reasoning Correlation}

\begin{table}[h]
\centering
\small
\caption{Correlation between retrieval and reasoning across all runs.
Values from \texttt{run\_correlation\_analysis.py}.}
\label{tab:correlation}
\begin{tabular}{lcccc}
\toprule
\textbf{Metric} & \textbf{Pearson $r$} & \textbf{$p$} & \textbf{Spearman $\rho$} & \textbf{$p$} \\
\midrule
Precision@K  & 0.244 & $<$0.001 & 0.222 & $<$0.001 \\
Recall@K     & 0.208 & $<$0.001 & 0.178 & 0.0004 \\
MRR          & 0.169 & 0.0008 & 0.143 & 0.0045 \\
\bottomrule
\end{tabular}
\end{table}

All correlations are weak ($|r|{<}0.25$, $|\rho|{<}0.23$), despite
being statistically significant ($p{<}0.005$) due to the large
sample ($n{=}390$).
Pearson $r{=}0.244$ for P@K means retrieval quality explains only
$\sim$6\% of variance in reasoning accuracy.
This confirms the gap: \textbf{retrieval quality is a poor predictor
of reasoning accuracy}.

\subsection{Failure Mode Distribution}

\begin{table}[h]
\centering
\small
\caption{Per-query regime classification.}
\label{tab:failure-modes}
\begin{tabular}{lcc}
\toprule
\textbf{Regime} & \textbf{Count} & \textbf{\%} \\
\midrule
High retrieval, high reasoning           & 111 & 28.5 \\
High retrieval, low reasoning            & 84 & 21.5 \\
Low retrieval, high reasoning            & 84 & 21.5 \\
Low retrieval, low reasoning             & 111 & 28.5 \\
\bottomrule
\end{tabular}
\end{table}

The 43\% of queries in off-diagonal regimes (21.5\% high-retrieval/low-reasoning + 21.5\% low-retrieval/high-reasoning)
are direct evidence that aggregate metrics obscure distinct failure
modes.
The symmetric distribution across quadrants further supports the
independence claim: retrieval success neither guarantees nor
precludes correct reasoning.

\subsection{Statistical Significance}

Table~\ref{tab:significance} reports paired $t$-tests (65 matched
queries) with bootstrap 95\% confidence intervals.

\begin{table}[h]
\centering
\small
\caption{Statistical significance of key comparisons (paired $t$-test, $n{=}65$).}
\label{tab:significance}
\begin{tabular}{llccc}
\toprule
\textbf{Comparison} & \textbf{Metric} & \textbf{$\Delta$} & \textbf{$p$} & \textbf{95\% CI} \\
\midrule
Hybrid $\to$ Hybrid+Rerank & P@K & +0.096 & $<$0.001*** & [0.057, 0.134] \\
Hybrid $\to$ Hybrid+Rerank & RCA & +0.032 & 0.234 & [$-$0.021, 0.085] \\
\addlinespace
Dense-Only $\to$ Full System & P@K & +0.198 & $<$0.001*** & [0.149, 0.247] \\
Dense-Only $\to$ Full System & RCA & +0.048 & 0.083 & [$-$0.006, 0.102] \\
\addlinespace
Hybrid $\to$ Full System & P@K & +0.053 & 0.009** & [0.013, 0.092] \\
Hybrid $\to$ Full System & RCA & +0.047 & 0.132 & [$-$0.014, 0.107] \\
\bottomrule
\end{tabular}
\end{table}

Retrieval improvements are consistently significant ($p{<}0.01$),
while reasoning differences are consistently non-significant
($p{>}0.08$).
This pattern---significant retrieval gains paired with
non-significant reasoning changes---is the statistical signature of
the retrieval--reasoning gap.

\subsection{Improved Metrics}
\label{sec:metrics-result}

\begin{table}[h]
\centering
\small
\caption{Metric inter-correlation.
Values from \texttt{run\_improved\_metrics.py}.}
\label{tab:metric-corr}
\begin{tabular}{lcc}
\toprule
\textbf{Metric Pair} & \textbf{Pearson $r$} \\
\midrule
Keyword overlap vs.\ Semantic similarity      & 0.673 \\
Keyword overlap vs.\ Composite score          & 0.773 \\
Semantic sim.\ vs.\ ROUGE-L F1                & 0.628 \\
\bottomrule
\end{tabular}
\end{table}

Keyword overlap correlates only moderately with semantic similarity
($r{=}0.673$), indicating it is a noisy proxy---approximately 55\%
of variance is unexplained.
The composite score correlates more strongly with semantic similarity
($r{=}0.849$), confirming that multi-signal evaluation reduces
single-metric vulnerability.

\subsection{Qualitative Cases}
\label{sec:cases}

\paragraph{Case 1: Good retrieval, wrong reasoning.}
On SCTP failure (S02), Hybrid+Rerank retrieves all relevant files
(P@K~${>}0.8$), but the LLM stops at the symptom without tracing the
full cascade.

\paragraph{Case 2: Imperfect retrieval, correct reasoning.}
On handover failure (S03), Dense-Only misses RAN entries, but the LLM
correctly chains the failure using domain priors.

\paragraph{Case 3: Iteration helps.}
On authentication failure (S07), first-pass retrieval is broad;
query refinement injects ``SQN'', ``AUSF'', pulling correct evidence
on iteration~2.

\paragraph{Case 4: Iteration hurts.}
Under Hybrid+Iter.\ on scheduler starvation (S08), incorrect
first-pass analysis causes query drift toward unrelated transport logs.

% ═══════════════════════════════════════════════════════════════════════════════
% 6. DISCUSSION
% ═══════════════════════════════════════════════════════════════════════════════

\section{Discussion}
\label{sec:discussion}

\subsection{Why Retrieval Improvements Fail to Improve Reasoning}

The gap has a structural explanation: retrieval and reasoning are
bounded by \emph{different failure modes}.
Retrieval is limited by index coverage, fusion weights, and reranker
accuracy.
Reasoning is limited by the LLM's ability to chain multi-source
evidence and distinguish relevant from irrelevant context.

Three mechanisms generate the disconnect:
(i)~Precision@K is reasoning-blind---a document can be relevant by
annotation yet unused by the LLM.
(ii)~End-task accuracy is retrieval-blind---correct answers can arise
from LLM priors rather than retrieved evidence.
(iii)~Averaging across queries erases the regime---two queries in
opposite failure modes produce identical aggregate scores.
The scatter plots in our supplementary materials (Figure~1) visualize
this: across 390 query$\times$config pairs, the regression slope is
near-flat ($r{=}0.244$), with large vertical spread confirming that
identical retrieval scores produce widely varying reasoning outcomes.

\subsection{Implications for RAG Evaluation}

Evaluating multi-hop RAG requires:
(i)~explicit decomposition of retrieval vs.\ reasoning failures;
(ii)~ablations holding the LLM fixed while varying retrieval, and
vice versa;
(iii)~per-query regime reporting;
(iv)~semantic evaluation metrics beyond keyword overlap.
These apply to multi-hop QA, scientific RAG, code assistants, and
agentic retrieval.

\subsection{Adaptive vs.\ Fixed Iteration}

Adaptive stopping avoids wasteful LLM calls and tolerates temporary
confidence dips.
Fixed iteration exhibits query drift (Case~4), amplifying first-pass
errors without a precision-oriented reranking guard.

% ═══════════════════════════════════════════════════════════════════════════════
% 7. LIMITATIONS AND FUTURE WORK
% ═══════════════════════════════════════════════════════════════════════════════

\section{Limitations and Future Work}
\label{sec:limitations}

\textbf{Synthetic data}: results should be validated on production logs.
\textbf{LLM-as-judge bias}: cross-model judges and human evaluation
are needed.
\textbf{Single domain}: replication on HotpotQA, MuSiQue, and
code-assistant workloads would test generality.
\textbf{Confidence calibration}: the composite score is a heuristic,
not a validated metric.

Future directions include: (i)~larger corpora with bootstrap
significance testing; (ii)~human-calibrated evaluation;
(iii)~domain-adapted embeddings via contrastive learning;
(iv)~graph-based causal retrieval; (v)~reader fine-tuning to test
whether reasoning improvements close the gap; (vi)~cross-domain
replication.

% ═══════════════════════════════════════════════════════════════════════════════
% 8. CONCLUSION
% ═══════════════════════════════════════════════════════════════════════════════

\section{Conclusion}

We demonstrated the \textbf{retrieval--reasoning gap}: across six
ablation configurations and 65 queries, P@K varies by 0.24
($p{<}0.001$) while RCA\textsubscript{kw} varies by only 0.05
($p{=}0.08$, n.s.), with Pearson $r{=}0.244$ between retrieval and
reasoning.
A multi-LLM comparison with fixed retrieval shows reasoning accuracy
gaps of up to 0.103 attributable to model capability alone.
Per-query analysis reveals 43\% of cases fall in off-diagonal
regimes where retrieval and reasoning quality diverge.
We proposed a decomposition framework and composite metric
($r{=}0.849$ with semantic evaluation) that make the gap visible.
The broader implication: \emph{for multi-hop RAG, improving retrieval
is necessary but not sufficient---reasoning quality requires
independent attention, and evaluation must reflect this.}

% ═══════════════════════════════════════════════════════════════════════════════
% REFERENCES
% ═══════════════════════════════════════════════════════════════════════════════

\bibliographystyle{plainnat}

\begin{thebibliography}{15}

\bibitem[Asai et~al.(2024)]{asai2024selfrag}
Asai, A., Wu, Z., Wang, Y., Sil, A., and Hajishirzi, H.
\newblock Self-RAG: Learning to retrieve, generate, and critique
through self-reflection.
\newblock In \emph{ICLR}, 2024.

\bibitem[Gao et~al.(2023)]{gao2023ragsurvey}
Gao, Y., Xiong, Y., Dibia, V., et~al.
\newblock Retrieval-augmented generation for large language models: A
survey.
\newblock \emph{arXiv preprint arXiv:2312.10997}, 2023.

\bibitem[Jiang et~al.(2023)]{jiang2023activerar}
Jiang, Z., Xu, F.~F., Gao, L., et~al.
\newblock Active retrieval augmented generation.
\newblock In \emph{EMNLP}, 2023.

\bibitem[Khattab and Zaharia(2020)]{khattab2020colbert}
Khattab, O. and Zaharia, M.
\newblock ColBERT: Efficient and effective passage search via
contextualized late interaction over BERT.
\newblock In \emph{SIGIR}, 2020.

\bibitem[Lewis et~al.(2020)]{lewis2020rag}
Lewis, P., Perez, E., Piktus, A., et~al.
\newblock Retrieval-augmented generation for knowledge-intensive NLP
tasks.
\newblock In \emph{NeurIPS}, 2020.

\bibitem[Nguyen et~al.(2016)]{nguyen2016msmarco}
Nguyen, T., Rosenberg, M., Song, X., et~al.
\newblock {MS MARCO}: A human generated machine reading comprehension
dataset.
\newblock In \emph{NeurIPS Workshop on Cognitive Computation}, 2016.

\bibitem[Nogueira and Cho(2019)]{nogueira2019passage}
Nogueira, R. and Cho, K.
\newblock Passage re-ranking with BERT.
\newblock \emph{arXiv preprint arXiv:1901.04085}, 2019.

\bibitem[Reimers and Gurevych(2019)]{reimers2019sentencebert}
Reimers, N. and Gurevych, I.
\newblock Sentence-BERT: Sentence embeddings using Siamese
BERT-networks.
\newblock In \emph{EMNLP}, 2019.

\bibitem[Robertson and Zaragoza(2009)]{robertson2009bm25}
Robertson, S. and Zaragoza, H.
\newblock The probabilistic relevance framework: {BM25} and beyond.
\newblock \emph{Foundations and Trends in IR}, 3(4):333--389, 2009.

\bibitem[Thakur et~al.(2021)]{thakur2021beir}
Thakur, N., Reimers, N., et~al.
\newblock {BEIR}: A heterogeneous benchmark for zero-shot evaluation
of information retrieval models.
\newblock In \emph{NeurIPS Datasets and Benchmarks}, 2021.

\bibitem[Wang et~al.(2023)]{wang2023queryrewrite}
Wang, L., Yang, N., and Wei, F.
\newblock Query2doc: Query expansion with large language models.
\newblock In \emph{EMNLP}, 2023.

\bibitem[Es et~al.(2024)]{es2024ragas}
Es, S., James, J., Espinosa-Anke, L., and Schockaert, S.
\newblock {RAGAS}: Automated evaluation of retrieval augmented
generation.
\newblock In \emph{EACL (System Demonstrations)}, 2024.

\bibitem[Saad-Falcon et~al.(2024)]{saad2024ares}
Saad-Falcon, J., Khattab, O., Potts, C., and Zaharia, M.
\newblock {ARES}: An automated evaluation framework for
retrieval-augmented generation systems.
\newblock In \emph{NAACL}, 2024.

\bibitem[Upadhyay et~al.(2024)]{upadhyay2024trecrag}
Upadhyay, S., Pradeep, R., Thakur, N., et~al.
\newblock A large-scale study of relevance assessments with large
language models: An initial look at the {TREC 2024 RAG Track}.
\newblock \emph{arXiv preprint arXiv:2411.08275}, 2024.

\bibitem[Liu et~al.(2024)]{liu2024lostmiddle}
Liu, N.~F., Lin, K., Hewitt, J., et~al.
\newblock Lost in the middle: How language models use long contexts.
\newblock \emph{TACL}, 12:157--173, 2024.

\end{thebibliography}

\end{document}
